\section{Introduction}

Python is among the most widely used languages in scientific computing, data analysis, and teaching, valued for
an accessible surface syntax and a large ecosystem. That reach comes with a large and permissive language.
Programs mutate state freely, depend on side effects and dynamic binding, and admit many implicit coercions, so
they can be hard to reason about, hard to reproduce, and hard to target with formal tools or with
implementations other than the reference interpreter. A researcher in programming languages or in pedagogy who
wants a well-behaved functional language faces a trade-off: such a language is easy to reason about but
unfamiliar, whereas Python is familiar but unruly.

\PurePy carves out a pure, side-effect-free subset of Python that is well-behaved and easy to reason about, yet
remains ordinary Python. Every well-formed \PurePy program is valid Python, and where \PurePy gives a run a
result, a standard interpreter gives the same one, so one artifact serves two roles: a program a scientist can run today, and a
common target that a research language or tool can produce and consume. Purity makes programs referentially
transparent, deterministic, and reproducible, the properties that scientific computing and formal reasoning
both want.

This paper defines \PurePy as a versioned formal standard\footnote{This paper describes version
\versionstring{} of the \PurePy language specification, which \specificationavailability.}: a formal grammar, a type system whose well-typed
programs mypy accepts, and an operational semantics that agrees with Python on every run it defines. A reference checker decides well-formedness, and a conformance test suite exercises the
standard against CPython,  GraalPy, and Pyodide. The guarantee is precise: well-formedness is decidable, and
\PurePy either gives a run Python's result or gives it no result at all.

The design is deliberately conservative. Rather than reconstruct Python's full behaviour, \PurePy chooses a
simple, well-behaved core and excludes the idioms that would break it, among them mutable state, effectful
control, implicit coercions, and constructs whose meaning depends on module load order
(\secref{overview}). The standard is versioned, so the subset can grow as features earn their place
without disturbing programs written against an earlier version.

\PurePy is aimed first at researchers in programming languages and in pedagogy, and is intended to grow into a
common language for scientific computing, supporting portable applications in modelling, data processing,
analysis, and visualisation. Because every \PurePy program is also a Python program, adopting it costs nothing
at the surface, and a research language that targets \PurePy inherits a large and familiar runtime.

\subsection{Contributions}

After setting out the problem \PurePy addresses in \secref{overview}, we make the following
contributions:
\begin{itemize}
\item a definition of \PurePy as a versioned standard, with a formal grammar, a type system
(\secref{well-formedness}) and an operational semantics that agrees with Python on
every run it defines (\secref{operational-semantics});
\item a reference checker that decides well-formedness, validated by a conformance test suite
exercised on CPython,  GraalPy, and Pyodide (\secref{reference-implementation});
\item conservative extensions to the core language, including partial application and piecewise function
definitions (\secref{extensions});
\item case studies that use \PurePy as a common functional language (\secref{case-studies});
\item an experimental evaluation of how well AI agents generate \PurePy and translate existing Python into it
(\secref{evaluation}).
\end{itemize}
\secref{related-work} summarises related efforts, and \secref{conclusion} presents our conclusions
and discusses future work.
